% META: {"id": "TCOMPL-ME-EV-B", "chapitre": "TCOMPL-MODELES-EVOLUTION", "type_objet": "evaluation", "version": "B", "duree_min": 50, "bareme_total": 20, "capacites_codes": ["C3", "C5"], "competences": ["calculer", "raisonner"], "mode_creation": "ex_nihilo", "sources_inspiration": [], "status": "generated"}
% BEGIN-VERIFY
% from sympy import symbols, exp, diff, Rational, limit, oo
% t = symbols('t')
% y = -80*exp(-Rational(1,3)*t) + 100
% yp = diff(y, t)
% assert (yp - (-Rational(1,3)*y + Rational(100,3))).simplify() == 0
% assert y.subs(t, 0) == 20
% assert limit(y, t, oo) == 100
% END-VERIFY

%---------------------------------------------------------------
% Duree : 50 minutes — Bareme : 20 points
%---------------------------------------------------------------

\begin{evaluation}{TCOMPL-ME-EV-B}{50}

\begin{exercice}{TCOMPL-ME-EV-B-EX1}{2}{20}
\textbf{Exercice 1} (8 points) --- \textit{Capacité C3}

\medskip

On définit $u_0=1$ et $u_{n+1} = \dfrac{1}{3}u_n + 4$.
\begin{enumerate}
  \item (3 pts) Un point fixe de $f(x)=\dfrac{1}{3}x+4$ correspond à l'intersection entre la courbe de $f$ et la droite $y=x$. Le déterminer par le calcul.
  \item (5 pts) Calculer $u_1$ et $u_2$, et comparer leur proximité avec le point fixe trouvé à la question précédente.
\end{enumerate}
\end{exercice}

\begin{exercice}{TCOMPL-ME-EV-B-EX2}{2}{30}
\textbf{Exercice 2} (12 points) --- \textit{Capacité C5}

\medskip

La quantité $y(t)$ d'un médicament dans le sang (en mg) vérifie $y'(t) = -\dfrac{1}{3}y(t) + \dfrac{100}{3}$, avec $y(0)=20$.
\begin{enumerate}
  \item (3 pts) Déterminer la solution particulière constante $y_p$.
  \item (5 pts) Donner la solution générale, puis déterminer $y(t)$ à l'aide de la condition initiale.
  \item (4 pts) Quelle est la quantité limite de médicament dans le sang lorsque $t\to+\infty$ ?
\end{enumerate}
\end{exercice}

\end{evaluation}
